\documentclass[12pt,a4paper]{article}

% -------------------------------------------------
% Мова та шрифти: компілювати через XeLaTeX або LuaLaTeX
% -------------------------------------------------
\usepackage{fontspec}
\usepackage{polyglossia}

\setmainlanguage{ukrainian}
\setotherlanguage{english}

\setmainfont{Libertinus Serif}
\setsansfont{Libertinus Sans}
\setmonofont{Libertinus Mono}

\usepackage{unicode-math}
\setmathfont{Libertinus Math}

% -------------------------------------------------
% Поля сторінки
% -------------------------------------------------
\usepackage[
a4paper,
left=25mm,
right=20mm,
top=20mm,
bottom=20mm
]{geometry}

% -------------------------------------------------
% Математика
% -------------------------------------------------
%\usepackage{amsmath}
%\usepackage{amssymb}
%\usepackage{mathtools}

% -------------------------------------------------
% Таблиці, рисунки, списки
% -------------------------------------------------
\usepackage{graphicx}
\usepackage{booktabs}
\usepackage{array}
\usepackage{enumitem}

\setlist[itemize]{leftmargin=1.2cm}
\setlist[enumerate]{leftmargin=1.2cm}

% -------------------------------------------------
% Абзаци та міжрядковий інтервал
% -------------------------------------------------
\usepackage{setspace}
\onehalfspacing

\setlength{\parindent}{1.25cm}
\setlength{\parskip}{0pt}

% -------------------------------------------------
% Посилання та URL
% -------------------------------------------------
\usepackage{url}
\usepackage[
unicode=true,
colorlinks=true,
linkcolor=black,
citecolor=black,
urlcolor=blue
]{hyperref}

% -------------------------------------------------
% Покращення перенесень та мікротипографіки
% -------------------------------------------------
\usepackage{microtype}
\sloppy

% -------------------------------------------------
% Назви розділів
% -------------------------------------------------
\usepackage{titlesec}

\titleformat{\section}
{\normalfont\Large\bfseries}
{\thesection}
{1em}
{}

\titleformat{\subsection}
{\normalfont\large\bfseries}
{\thesubsection}
{1em}
{}


\title{Механіка руйнування метаматеріально-орієнтованих інфраструктурних композитів з акцентом на системи дорожнього покриття за статичного та динамічного навантаження}

\author{Інститут механіки ім. С.П. Тимошенка НАН України}
\date{}

\begin{document}

\maketitle

\section*{Ключові слова}
% Для пункту 13 Запиту.
механіка руйнування; метаматеріально-орієнтовані композити; інфраструктурні композити; дорожнє покриття; метабетон; ауксетичні структури; когезійна зона; міжфазне руйнування; динамічне навантаження; тріщиностійкість.

% ------------------------------------------------------------
\section{Резюме}
% Для пункту 14 Запиту.
Сучасні інфраструктурні об'єкти, зокрема системи дорожнього покриття, мостові настили, плити, захисні та ремонтні шари, працюють в умовах складної дії статичних, циклічних і динамічних навантажень, температурних змін, вологості та поступового накопичення пошкоджень. Одним із головних чинників втрати їхньої довговічності є зародження і поширення тріщин у матеріалі або на міжфазних межах між його структурними компонентами. У зв'язку з цим актуальною є розробка механіко-математичних підходів до прогнозування тріщиностійкості, керування траєкторіями тріщин та підвищення опору руйнуванню композитних матеріалів інфраструктурного призначення.

У роботі пропонується дослідити метаматеріально-орієнтовані інфраструктурні композити, у яких ефективні механічні властивості формуються не лише складом матеріалу, а й спеціально спроєктованою внутрішньою архітектурою. Особливу увагу буде приділено композитам для систем дорожнього покриття, а також спорідненим інфраструктурним елементам, у яких внутрішні комірчасті, ауксетичні, резонансні або армувальні структури можуть впливати на розподіл напружень, локалізацію деформацій, дисипацію енергії та розвиток тріщин. Планується розроблення скінченноелементних і когезійно-зонних моделей зародження, поширення, відхилення та зупинки тріщин у таких матеріалах за статичного й динамічного навантаження, з урахуванням міжфазної деградації та можливої взаємодії тріщин із архітектурними елементами композиту.

Очікуваними результатами роботи є механіко-математичні моделі руйнування метаматеріально-орієнтованих інфраструктурних композитів, алгоритми визначення умов зародження та розвитку тріщин, чисельні карти руйнування для характерних внутрішніх архітектур, методики оцінювання міжфазної тріщиностійкості та рекомендації щодо вибору геометричних і матеріальних параметрів композитів із підвищеним опором пошкодженню. Результати дослідження можуть бути використані при створенні нових композитних матеріалів для дорожньої та транспортної інфраструктури, захисних і ремонтних шарів, плитних елементів, мостових настилів та інших інженерних об'єктів, для яких критичними є тріщиностійкість, довговічність і здатність до роботи за складних умов навантаження.


% ------------------------------------------------------------
\section{Вступ та актуальність дослідження}
% Основа для пункту 15.1 Запиту.

Дорожня, транспортна та інша цивільна інфраструктура працює в умовах тривалої дії статичних, циклічних і динамічних навантажень, температурних змін, вологи, локального зношування та поступового накопичення пошкоджень. Для дорожніх покриттів, мостових настилів, плитних елементів, захисних шарів і ремонтних прошарків одним із головних механізмів втрати працездатності є зародження та поширення тріщин. Тріщини можуть виникати як у цементній, бетонній, асфальтобетонній або полімерцементній матриці, так і на міжфазних межах між окремими структурними компонентами матеріалу. Їхній розвиток призводить до зниження жорсткості, несучої здатності, водонепроникності, морозостійкості та загальної довговічності інфраструктурних об'єктів.

Традиційні підходи до підвищення довговічності інфраструктурних композитів зазвичай пов'язані з удосконаленням складу матеріалу: використанням мінеральних і полімерних модифікаторів, волокнистого армування, високоміцних цементних матриць, спеціальних заповнювачів або ремонтних сумішей. Такі підходи залишаються важливими, однак вони не завжди забезпечують достатній контроль над траєкторіями поширення тріщин, локалізацією пошкоджень, міжфазним руйнуванням і динамічною відповіддю матеріалу. У зв'язку з цим актуальним стає перехід від переважно композиційного проєктування матеріалу до поєднання композиційного та архітектурного проєктування, коли внутрішня геометрія, топологія та зв'язність структурних елементів розглядаються як самостійні параметри керування механічною поведінкою.

Перспективним напрямом такого підходу є створення метаматеріально-орієнтованих інфраструктурних композитів. У цих матеріалах ефективні властивості визначаються не лише фізико-механічними характеристиками окремих фаз, а й спеціально спроєктованою внутрішньою архітектурою. Комірчасті структури, ауксетичні елементи, резонансні включення, перфоровані цементні компоненти, 3D-друковані армувальні решітки або полімерні вставки можуть змінювати шляхи передавання навантаження, перерозподіляти напруження, підвищувати енергопоглинання, формувати віброзахисні властивості та впливати на розвиток локальних пошкоджень \cite{Moini2024,Premadasa2026,Wang2025,Alavi2026}. Саме тому ідеї механічних метаматеріалів поступово переходять із лабораторних демонстрацій у сферу матеріалів для цивільної інфраструктури.

Для дорожніх та інших інфраструктурних систем особливо важливо, що внутрішня архітектура композиту може потенційно використовуватися не лише для зміни жорсткості, коефіцієнта Пуассона, хвильової відповіді чи енергопоглинання, а й для керування руйнуванням. Архітектурні елементи можуть виконувати функцію бар'єрів для тріщин, змінювати напрям їхнього поширення, сприяти дисипації енергії або забезпечувати стабільне посттріщинне деформування. Разом з тим ці самі елементи можуть створювати концентратори напружень, слабкі міжфазні межі або зони передчасного руйнування. Тому для таких матеріалів недостатньо оцінювати лише ефективну жорсткість, міцність або динамічну функціональність; необхідним є спеціальний аналіз з позицій механіки руйнування.

Сучасні дослідження вже демонструють окремі приклади ефективного використання метаматеріальних принципів у бетонних і цементних системах. Концепція метабетону з резонансними включеннями показує можливість керування динамічною реакцією бетонного композиту \cite{Mitchell2014}. Ауксетичні цементні та бетонні структури демонструють підвищене енергопоглинання, незвичайну деформаційну поведінку та роль crack-bridging механізмів у стабілізації посттріщинного деформування \cite{Xu2020,Xie2025,Chen2024,Vo2025}. Бетони з вбудованим 3D-друкованим архітектурним армуванням відкривають можливості для перерозподілу напружень і керування траєкторіями тріщин \cite{Vandadi2025}. Паралельно розвиваються багатофункціональні та самодіагностичні бетонні системи, орієнтовані на smart infrastructure, structural health monitoring та дорожні застосування \cite{Dong2023,Qin2024,Barri2023}.

Попри ці досягнення, системна механіка руйнування метаматеріально-орієнтованих інфраструктурних композитів залишається недостатньо розробленою. Більшість наявних робіт зосереджена на демонстрації окремих ефектів: ауксетичної поведінки, хвильового гасіння, енергопоглинання, сенсорної функціональності або підвищення міцності. Водночас недостатньо дослідженими залишаються закономірності зародження, поширення, відхилення та зупинки тріщин у таких матеріалах, вплив внутрішньої архітектури на crack-driving force, роль міжфазної деградації, а також довговічність полімерних або полімер-композитних елементів у цементному, лужному чи вологому середовищі \cite{Zhang2025,Jipa2022}.

Актуальність запропонованого дослідження полягає у необхідності створення науково обґрунтованих моделей і методик, які дадуть змогу перейти від окремих прикладів метаматеріальних ефектів до цілеспрямованого проєктування тріщиностійких інфраструктурних композитів. Особливий акцент робиться на системах дорожнього покриття, однак отримані результати можуть бути поширені також на мостові настили, плитні елементи, захисні та ремонтні шари, віброзахисні елементи й інші об'єкти цивільної інфраструктури. Розроблення таких підходів має наукове значення для розвитку механіки руйнування композитних матеріалів з внутрішньою архітектурою та практичне значення для підвищення довговічності, надійності й експлуатаційної стійкості інфраструктурних систем.


% ------------------------------------------------------------
\section{Наукова ідея та робоча гіпотеза}

Наукова ідея роботи полягає у використанні внутрішньої архітектури метаматеріально-орієнтованих інфраструктурних композитів як активного чинника керування руйнуванням. На відміну від традиційного підходу, за якого підвищення довговічності композиту досягається переважно зміною складу матеріалу, у цій роботі внутрішня геометрія, топологія комірок, розміщення включень, характер міжфазних меж і зв'язність архітектурних елементів розглядаються як параметри, що можуть цілеспрямовано впливати на зародження, поширення, відхилення та зупинку тріщин.

Основна ідея полягає в тому, що метаматеріальна архітектура може бути використана не лише для керування ефективною жорсткістю, коефіцієнтом Пуассона, енергопоглинанням або динамічною відповіддю матеріалу, а й для формування внутрішніх механізмів підвищення тріщиностійкості. Такими механізмами можуть бути перерозподіл напружень, зменшення локальної crack-driving force, відхилення траєкторії тріщини, її кiнкування на міжфазних межах, стабілізація посттріщинного деформування завдяки crack-bridging ефектам, а також створення архітектурних елементів, які виконують функцію внутрішніх ``пасток тріщин''.

Робоча гіпотеза дослідження полягає в тому, що для інфраструктурних композитів з внутрішньою комірчастою, ауксетичною, резонансною, перфорованою або армувальною архітектурою можна підібрати такі геометричні, матеріальні та міжфазні параметри, за яких внутрішня структура не прискорює локальне руйнування, а навпаки підвищує опір зародженню і поширенню тріщин. Передбачається, що за належного поєднання архітектури композиту, посттріщинної поведінки матриці та властивостей міжфазних меж можна досягти підвищеного енергопоглинання, стабільнішого розвитку пошкоджень і більшої довговічності матеріалу за статичного, циклічного та динамічного навантаження.

Об'єктом дослідження є метаматеріально-орієнтовані інфраструктурні композити, зокрема композити для систем дорожнього покриття, плитних елементів, мостових настилів, захисних і ремонтних шарів, у яких внутрішня архітектура істотно впливає на напружено-деформований стан, локалізацію пошкоджень і розвиток тріщин.

Предметом дослідження є закономірності зародження, поширення, відхилення та зупинки тріщин у метаматеріально-орієнтованих інфраструктурних композитах за статичного й динамічного навантаження, а також вплив внутрішньої архітектури, міжфазної деградації, crack-bridging механізмів і середовищних чинників на тріщиностійкість і довговічність таких матеріалів.

Перевірка робочої гіпотези передбачає побудову скінченноелементних і когезійно-зонних моделей руйнування, аналіз взаємодії тріщин з архітектурними елементами композиту, дослідження міжфазного руйнування та чисельно-експериментальну верифікацію отриманих результатів. Такий підхід дасть змогу перейти від опису окремих метаматеріальних ефектів до цілеспрямованого проєктування інфраструктурних композитів, у яких внутрішня архітектура використовується як інструмент підвищення опору руйнуванню.


\subsection{Об'єкт дослідження}

Об'єктом дослідження є метаматеріально-орієнтовані інфраструктурні композити з внутрішньою архітектурою, що впливає на їхню механічну поведінку, тріщиностійкість і довговічність. До таких композитів належать, зокрема, матеріали для систем дорожнього покриття, плитних елементів, мостових настилів, захисних і ремонтних шарів, у яких комірчасті, ауксетичні, резонансні, перфоровані, армувальні або міжфазні структурні елементи можуть змінювати розподіл напружень, локалізацію деформацій і розвиток пошкоджень.


\subsection{Предмет дослідження}

Предметом дослідження є закономірності зародження, поширення, відхилення та зупинки тріщин у метаматеріально-орієнтованих інфраструктурних композитах за статичного, циклічного та динамічного навантаження, а також вплив внутрішньої архітектури, міжфазної деградації, crack-bridging механізмів, посттріщинної поведінки матриці та середовищних чинників на тріщиностійкість, енергопоглинання і довговічність таких матеріалів.


% ------------------------------------------------------------
\section{Стан розроблення проблеми}

У сучасних дослідженнях механічні метаматеріали розглядаються як перспективний клас матеріалів і структур для цивільної інфраструктури, оскільки їхні ефективні властивості визначаються не лише складом, а й спеціально спроєктованою внутрішньою архітектурою. Такий підхід дає змогу керувати жорсткістю, енергопоглинанням, деформаційною здатністю, віброізоляційними властивостями, адаптивністю та сенсорними функціями матеріалу \cite{Moini2024,Premadasa2026,Wang2025,Alavi2026}. На відміну від традиційних бетонів, цементних композитів або полімерцементних систем, у метаматеріально-орієнтованих композитах внутрішня геометрія, зв'язність елементів, характер міжфазних меж і топологія шляхів передавання навантаження можуть розглядатися як самостійні параметри проєктування. Це відкриває можливості для створення нових матеріалів інфраструктурного призначення, зокрема для дорожніх покриттів, плит, мостових настилів, захисних шарів, ремонтних прошарків і віброзахисних елементів.

Важливим напрямом розвитку цієї галузі є архітектурні цементні та бетонні матеріали. У таких системах механічна поведінка квазікрихкої матриці може бути істотно змінена завдяки введенню регулярної або спеціально оптимізованої внутрішньої структури. У роботах, присвячених architected infrastructure materials, підкреслюється, що для цементних і бетонних матеріалів внутрішня архітектура може стати не лише засобом підвищення міцності або зменшення маси, а й інструментом керування тріщиностійкістю, дисипацією енергії, локалізацією пошкоджень і стабільністю посттріщинного деформування \cite{Moini2024}. Особливо важливим є те, що для таких матеріалів проблема руйнування не може розглядатися як другорядна: саме механізми зародження тріщин, їхнього поширення, відхилення або зупинки визначають придатність архітектурних композитів до роботи в реальних інфраструктурних умовах.

Одним із перших прикладів перенесення ідей метаматеріалів у бетонні композити є концепція метабетону, у якій звичайні заповнювачі замінюються спеціально спроєктованими резонансними включеннями. Такі включення можуть складатися з важкого ядра та м'якого покриття, що утворює локальний резонатор, здатний поглинати частину енергії динамічного навантаження та зменшувати рівень напружень у крихкій матриці \cite{Mitchell2014}. Цей підхід особливо важливий для задач ударного, імпульсного, вибухового та вібраційного навантаження інфраструктурних елементів. Разом з тим у наявних роботах основну увагу приділено хвильовій відповіді та ефектам локального резонансу, тоді як питання впливу резонансних включень на зародження тріщин, crack-driving force, міжфазне відшарування та динамічне поширення тріщини залишаються недостатньо розробленими.

Близьким за змістом напрямом є цементні метаматеріали для віброізоляції та хвильового керування. Зокрема, запропоновано цементні структури зі спіральними або щілинними резонаторами, які можуть формувати смуги заборонених частот і зменшувати передавання низькочастотних коливань \cite{Imagawa2025}. Такі рішення демонструють принципову можливість створення однофазних цементних метаматеріалів без введення додаткових важких або м'яких включень. Однак з погляду механіки руйнування щілинні та резонаторні елементи є також концентраторами напружень, у вершинах яких можуть зароджуватися тріщини за статичного, циклічного або динамічного навантаження. Тому для інфраструктурних застосувань необхідно пов'язувати динамічну функціональність таких структур з оцінкою їхньої тріщиностійкості та довговічності.

Значний розвиток отримали ауксетичні цементні та бетонні композити. У цих матеріалах внутрішня комірчаста або перфорована архітектура забезпечує незвичайну деформаційну поведінку, зокрема від'ємний або близький до нуля коефіцієнт Пуассона, підвищене енергопоглинання та здатність до більш стабільного посттріщинного деформування. Показано, що у цементних комірчастих композитах ауксетична поведінка може виникати не лише внаслідок пружного згинання елементів, а й завдяки crack-bridging механізму: тріщини зароджуються у шарнірних зонах комірок, але місткове зчеплення волокон утримує береги тріщин і дає змогу елементам обертатися без катастрофічного розділення \cite{Xu2020}. Подальші дослідження показали, що використання strain-hardening cementitious composites істотно покращує несучу здатність, стабільність ауксетичного деформування, енергопоглинання та відновлювану деформацію завдяки затримці локалізації магістральної тріщини і зменшенню її розкриття \cite{Xie2025}.

Порівняльні дослідження різних ауксетичних архітектур на основі engineered cementitious composites свідчать, що топологія комірки істотно впливає на характер руйнування, жорсткість, міцність, згинальну поведінку та здатність до дисипації енергії. Зокрема, re-entrant, chiral та buckling-induced архітектури демонструють різні механізми локалізації пошкоджень, а buckling-induced структури можуть забезпечувати кращу несучу здатність і вищу згинальну тріщиностійкість \cite{Chen2024}. У дослідженнях ауксетичного метабетону на основі UHPFRC також показано, що ефективність матеріалу визначається одночасно геометрією перфорацій і посттріщинною поведінкою матриці. Зокрема, strain-hardening UHPFRC із розвиненим містковим зчепленням волокон забезпечує більш стабільне обертання комірок, збереження структурної цілісності та підвищене енергопоглинання \cite{Vo2025}. Отже, для таких матеріалів руйнування не є лише небажаним граничним станом, а стає частиною керованого механізму деформування, який повинен описуватися методами механіки руйнування.

Інший перспективний підхід полягає у введенні в бетонну або цементну матрицю 3D-друкованих ауксетичних чи метаматеріальних армувальних елементів. На відміну від випадкового дисперсного армування, така внутрішня структура може цілеспрямовано перерозподіляти напруження, змінювати траєкторію тріщини, затримувати її поширення та підвищувати енергопоглинання композиту \cite{Vandadi2025}. У таких системах особливого значення набуває взаємодія тріщини з армувальним елементом: тріщина може проникати в армування, відхилятися вздовж межі поділу, кiнкувати в матрицю або зупинятися на архітектурному бар'єрі. Це створює передумови для проєктування внутрішніх ``пасток тріщин'', але водночас потребує спеціальних моделей міжфазного руйнування та контактної взаємодії.

Паралельно розвиваються багатофункціональні та самодіагностичні метаматеріальні бетони. Сучасні дослідження smart pavement concretes та electricity-based multifunctional concrete демонструють можливості самодіагностики, моніторингу напружено-деформованого стану, енергозбирання, протиожеледних функцій, traffic detection та structural health monitoring \cite{Dong2023,Qin2024}. Окремий приклад становлять nanogenerator-integrated metamaterial concrete systems, у яких ауксетична полімерна решітка з snap-through поведінкою поєднується з провідною цементною матрицею та трибоелектричним механізмом енергозбирання \cite{Barri2023}. Такі матеріали можуть мати самовідновлюване деформування, сенсорну функціональність і здатність до виявлення пошкоджень. Водночас для запропонованого проєкту ці роботи мають передусім допоміжне значення: вони показують перспективу багатофункціональності, але не замінюють необхідності побудови механіки зародження, розвитку й накопичення тріщин.

Особливо важливою для інфраструктурних метакомпозитів є проблема міжфазного та середовищно-індукованого руйнування. У ремонтних і гібридних дорожніх системах полімер-бетонні або цемент-полімерні межі можуть бути зонами зародження крихкого руйнування, особливо за низьких температур, зсувного навантаження або неповного заповнення тріщин \cite{Zhang2025}. Когезійні моделі та енергетичні критерії змішаного руйнування є природним інструментом для опису таких процесів. Додатковим обмеженням є довговічність 3D-друкованих полімерів у контакті з цементним середовищем. Показано, що деякі полімери, зокрема PLA та його модифікації, можуть зазнавати швидкого environmental stress cracking за сумісної дії лужного середовища та розтягувальних напружень, тоді як ABS, ASA, HIPS, Nylon та деякі інші матеріали демонструють вищу стійкість \cite{Jipa2022}. Це означає, що вибір полімерної фази, орієнтація 3D-друку, стан міжфазної межі та умови експлуатаційного середовища повинні бути включені до fracture-mechanics аналізу бетонополімерних і цементополімерних метакомпозитів.

Таким чином, сучасна література підтверджує значний потенціал метаматеріально-орієнтованих інфраструктурних композитів для керування жорсткістю, енергопоглинанням, хвильовою відповіддю, деформаційною здатністю, сенсорною функціональністю та довговічністю. Водночас наявні дослідження переважно зосереджені на окремих ефектах: резонансному гасінні коливань, ауксетичному деформуванні, енергозбиранні, self-sensing або підвищенні міцності. Недостатньо розробленою залишається системна механіка руйнування таких матеріалів, яка б одночасно враховувала архітектуру комірок, crack-bridging поведінку матриці, міжфазну деградацію, динамічне навантаження та вплив експлуатаційного середовища. Саме ця наукова прогалина визначає актуальність запропонованого дослідження, спрямованого на побудову механіко-математичних моделей руйнування метаматеріально-орієнтованих інфраструктурних композитів з акцентом на системи дорожнього покриття.

% ------------------------------------------------------------
\section{Критичні наукові прогалини}

Аналіз сучасного стану досліджень показує, що метаматеріально-орієнтовані інфраструктурні композити розвиваються у кількох взаємопов'язаних напрямах: метабетон із резонансними включеннями, цементні метаматеріали для віброізоляції, ауксетичні цементні та бетонні структури, бетон із вбудованим архітектурним армуванням, а також багатофункціональні та самодіагностичні бетонополімерні системи. Однак більшість наявних робіт зосереджена на окремих ефектах --- енергопоглинанні, зменшенні маси, віброізоляції, ауксетичній деформаційній поведінці, сенсорних функціях або динамічному гасінні коливань. Значно менш розробленою залишається системна механіка руйнування таких матеріалів, особливо з урахуванням їхнього використання в дорожніх та інших інфраструктурних системах.

Першою критичною прогалиною є відсутність узагальненої механіко-математичної моделі руйнування метаматеріально-орієнтованих інфраструктурних композитів. У наявних дослідженнях внутрішня архітектура матеріалу зазвичай розглядається як засіб керування жорсткістю, ефективним коефіцієнтом Пуассона, хвильовою відповіддю або енергопоглинанням. Водночас недостатньо вивчено, як ця архітектура впливає на зародження тріщин, їхнє поширення, зупинку, розгалуження або відхилення траєкторії.

Другою прогалиною є недостатній розвиток моделей взаємодії тріщин з архітектурними елементами композиту. Комірчасті структури, ауксетичні вузли, перфорації, щілинні резонатори, армувальні решітки та включення можуть одночасно зменшувати локальні напруження і створювати нові концентратори напружень. Тому для таких матеріалів необхідно встановити умови, за яких внутрішня архітектура працює як механізм підвищення тріщиностійкості, а не як джерело передчасного руйнування.

Третьою прогалиною є недостатнє врахування посттріщинної поведінки цементної або бетонної матриці. У багатьох ауксетичних цементних композитах бажана деформаційна поведінка реалізується лише тоді, коли тріщини в шарнірних або лігаментних зонах не призводять до катастрофічного розділення елементів, а стабілізуються завдяки crack-bridging механізмам, волокнистому армуванню або strain-hardening поведінці матриці. Отже, для прогнозування механічної відповіді таких систем необхідно поєднувати геометричний опис архітектури з когезійними або енергетичними моделями посттріщинного деформування.

Четвертою прогалиною є недостатній зв'язок між динамічною функціональністю метаматеріальних систем і критеріями руйнування. Метабетон, резонансні включення та цементні метаматеріали зі щілинними або спіральними резонаторами здатні змінювати хвильову відповідь і зменшувати передавання динамічної енергії. Проте залишається відкритим питання, як ці ефекти впливають на локальні параметри руйнування: максимальні розтягувальні напруження, енергію руйнування, коефіцієнти інтенсивності напружень, crack-driving force та умови динамічного поширення тріщин.

П'ятою прогалиною є недостатньо розроблене моделювання міжфазного руйнування у гібридних бетонополімерних і цементополімерних метакомпозитах. У таких системах межі полімер--цемент, армування--матриця, включення--покриття--матриця або ремонтний шар--основа можуть бути зонами зародження пошкодження. Тріщина може поширюватися вздовж інтерфейсу, проникати в матрицю, відхилятися на межі поділу або спричиняти локальне відшарування архітектурного елемента. Для опису цих процесів необхідні когезійні моделі змішаного руйнування та критерії переходу тріщини між фазами композиту.

Шостою прогалиною є недостатнє врахування довговічності полімерних і полімер-композитних архітектур у цементному, лужному, вологому або температурно змінному середовищі. Для 3D-друкованих полімерних решіток, форм, вставок або армувальних елементів механічна працездатність визначається не лише початковою міцністю матеріалу, а й його хіміко-механічною сумісністю з цементною матрицею, орієнтацією друку, станом міжшарових меж і схильністю до середовищно-індукованого розтріскування. Це особливо важливо для інфраструктурних об'єктів, що працюють у довготривалих експлуатаційних умовах.

Таким чином, ключова наукова прогалина полягає у відсутності цілісного fracture-mechanics підходу до аналізу та проєктування метаматеріально-орієнтованих інфраструктурних композитів. Необхідно розробити моделі, які одночасно враховують внутрішню архітектуру, механізми зародження і поширення тріщин, посттріщинне деформування, міжфазну деградацію, динамічне навантаження та вплив експлуатаційного середовища. Заповнення цієї прогалини дасть змогу перейти від окремих демонстрацій метаматеріальних ефектів до науково обґрунтованого проєктування тріщиностійких композитів для дорожньої та іншої цивільної інфраструктури.


\begin{enumerate}
	\item Недостатньо розроблена системна механіка руйнування метаматеріально-орієнтованих інфраструктурних композитів.
	\item Недостатньо вивчена взаємодія тріщин з архітектурними елементами: порами, лігаментами, резонаторами, ауксетичними вузлами та вбудованим армуванням.
	\item Потребує розвитку моделювання міжфазного руйнування на межах полімер--цемент, армування--матриця, покриття--матриця та soft layer--mortar.
	\item Динамічна функціональність метаматеріальних систем часто не пов'язана з оцінкою тріщиностійкості та crack-driving force.
	\item Недостатньо враховано довговічність полімерних або полімер-композитних архітектур у цементному, лужному та вологому середовищі.
\end{enumerate}

% ------------------------------------------------------------
\section{Мета і завдання роботи}
% Для пункту 15.1 Запиту і пункту 6 Технічного завдання.

\subsection{Мета роботи}

Метою роботи є розроблення механіко-математичних моделей, чисельних методик та критеріїв оцінювання руйнування метаматеріально-орієнтованих інфраструктурних композитів з акцентом на системи дорожнього покриття за статичного та динамічного навантаження, з урахуванням впливу внутрішньої архітектури, міжфазної деградації, crack-bridging механізмів та взаємодії тріщин з архітектурними елементами матеріалу.

\subsection{Основні завдання роботи}

Для досягнення поставленої мети передбачається розв'язання таких основних завдань:

\begin{enumerate}
	\item Провести аналіз сучасного стану досліджень метаматеріально-орієнтованих інфраструктурних композитів, зокрема метабетону, ауксетичних цементних і бетонних структур, бетонів із 3D-друкованим архітектурним армуванням, багатофункціональних дорожніх матеріалів та бетонополімерних композитів.
	
	\item Визначити характерні типи внутрішньої архітектури інфраструктурних композитів, перспективні з погляду керування руйнуванням: комірчасті та ауксетичні структури, резонансні включення, перфоровані цементні елементи, полімерні або полімер-композитні армувальні решітки, а також шаруваті й міжфазні системи.
	
	\item Розробити скінченноелементні моделі напружено-деформованого стану метаматеріально-орієнтованих інфраструктурних композитів за статичного навантаження з урахуванням впливу геометричних параметрів внутрішньої архітектури на концентрацію напружень, локалізацію деформацій і умови зародження тріщин.
	
	\item Побудувати когезійно-зонні моделі руйнування характерних елементів архітектурних композитів, зокрема шарнірних зон ауксетичних комірок, лігаментів, меж включення--матриця, полімер--цементних і цемент--полімерних інтерфейсів, з урахуванням crack-bridging механізмів та посттріщинної поведінки матеріалу.
	
	\item Дослідити взаємодію тріщин із внутрішньою архітектурою композиту, включаючи відхилення траєкторії тріщини, її кiнкування, зупинку, проходження через архітектурний елемент або поширення вздовж міжфазної межі, та сформулювати критерії ефективності внутрішніх ``пасток тріщин''.
	
	\item Розробити чисельні моделі динамічного навантаження метаматеріально-орієнтованих композитів, зокрема композитів із резонансними включеннями або віброзахисними архітектурними елементами, та оцінити вплив динамічної функціональності на напружений стан, дисипацію енергії та параметри руйнування.
	
	\item Проаналізувати вплив міжфазної деградації, властивостей полімерних або полімер-композитних елементів, а також можливого середовищно-індукованого пошкодження на тріщиностійкість і довговічність бетонополімерних та цементополімерних метакомпозитів.
	
	\item Розробити підходи до експериментальної верифікації запропонованих моделей із використанням випробувань на стиск, згин, непрямий розтяг та SCB-випробувань для визначення параметрів тріщиностійкості, енергопоглинання, посттріщинної поведінки та міжфазного руйнування.
	
	\item Побудувати чисельні карти руйнування для вибраних типів метаматеріально-орієнтованих інфраструктурних композитів та сформулювати рекомендації щодо вибору геометричних, матеріальних і міжфазних параметрів, які забезпечують підвищений опір зародженню й поширенню тріщин.
	
	\item Обґрунтувати можливі напрями застосування отриманих результатів для систем дорожнього покриття, мостових настилів, плитних елементів, захисних і ремонтних шарів та інших інфраструктурних об'єктів, що працюють за статичного, циклічного або динамічного навантаження.
\end{enumerate}


% ------------------------------------------------------------
\section{Методи дослідження}
% Для пунктів 15.1, 15.4 і Технічного завдання.

\begin{itemize}
	\item метод скінченних елементів;
	\item когезійно-зонне моделювання;
	\item енергетичні критерії руйнування;
	\item моделі міжфазної деградації;
	\item моделювання статичного, циклічного та імпульсного динамічного навантаження;
	\item параметричний аналіз геометрії комірок і включень;
	\item експериментальна верифікація за схемами стиску, згину, непрямого розтягу та SCB.
\end{itemize}

% ------------------------------------------------------------
\section{Заплановані чисельні симуляції}

Чисельна частина дослідження буде спрямована на встановлення зв'язку між внутрішньою архітектурою метаматеріально-орієнтованих інфраструктурних композитів і механізмами їхнього руйнування. Основну увагу буде приділено моделюванню зародження, поширення, відхилення та зупинки тріщин, а також оцінюванню ролі міжфазних меж, crack-bridging механізмів, динамічного навантаження і середовищної деградації. Для цього передбачається виконання кількох взаємопов'язаних груп чисельних симуляцій.

\subsection*{Симуляція 1. Ауксетична цементна комірка з когезійними шарнірними зонами}

Буде розроблено скінченноелементну модель ауксетичної цементної або бетонної комірки, у якій деформаційна поведінка визначається обертанням лігаментів і локальним тріщиноутворенням у шарнірних зонах. Такі моделі відповідають сучасним підходам до опису cementitious cellular composites та auxetic cementitious materials, у яких тріщини в зонах з'єднання елементів можуть не призводити до негайного руйнування, а стабілізуватися завдяки crack-bridging механізмам \cite{Xu2020,Xie2025}.

У моделі передбачається введення когезійних елементів або когезійних законів у потенційних зонах зародження тріщин. Буде досліджено вплив параметрів когезійного закону, жорсткості шарнірної зони, міцності на розтяг, енергії руйнування та геометрії комірки на характер деформування, apparent Poisson's ratio, локалізацію пошкоджень, енергопоглинання та залишкову несучу здатність. Окремо буде проаналізовано перехід від крихкого руйнування шарнірної зони до стабільного посттріщинного деформування за наявності ефекту місткового зчеплення.

Очікуваний результат цієї групи симуляцій --- визначення умов, за яких локальне тріщиноутворення в ауксетичній комірці стає не катастрофічним механізмом руйнування, а контрольованим механізмом деформування й дисипації енергії.

\subsection*{Симуляція 2. Порівняння різних архітектур цементних і бетонних метакомпозитів за стиску та згину}

Буде виконано порівняльне моделювання кількох типів внутрішньої архітектури цементних і бетонних композитів: re-entrant, chiral, buckling-induced, diamond, elliptical та peanut-like структур. Такі архітектури розглядаються у сучасних дослідженнях ECC- та UHPFRC-based auxetic structures як перспективні для підвищення енергопоглинання, несучої здатності та стабільності посттріщинного деформування \cite{Chen2024,Vo2025}.

Для кожної архітектури буде побудовано моделі навантаження за стиску та згину. Передбачається аналіз розподілу головних розтягувальних напружень, зон зародження тріщин, траєкторій їхнього поширення, характеру локалізації пошкоджень і впливу відносної густини архітектури на механічну відповідь. У разі необхідності в критичних зонах буде введено когезійні елементи для опису ініціювання і розвитку тріщин.

Очікуваний результат --- побудова порівняльних чисельних карт руйнування для різних типів архітектури та визначення геометричних параметрів, які забезпечують найкраще поєднання жорсткості, тріщиностійкості, енергопоглинання і стабільного посттріщинного деформування.

\subsection*{Симуляція 3. Метабетон із резонансними включеннями за імпульсного або динамічного навантаження}

Буде розроблено модель бетонного або цементного композиту з резонансними включеннями, що складаються з ядра та м'якого покриття. Така схема відповідає концепції метабетону, у якій локальні резонатори здатні перерозподіляти механічну енергію і зменшувати інтенсивність динамічної дії на крихку матрицю \cite{Mitchell2014}. Також буде враховано ідеї цементних метаматеріалів із щілинними або спіральними резонаторними елементами \cite{Imagawa2025}.

У симуляціях передбачається аналіз проходження імпульсного навантаження або хвильового збурення крізь модельний об'єм композиту. Буде досліджено вплив розміру, жорсткості, густини, розміщення і частотного налаштування включень на розподіл енергії, максимальні розтягувальні напруження у матриці, локалізацію пошкоджень і параметри руйнування. Окремо буде проаналізовано можливість зародження тріщин біля меж включення--матриця та роль когезійного руйнування м'якого покриття або міжфазної межі.

Очікуваний результат --- встановлення зв'язку між динамічною функціональністю метабетону та параметрами руйнування, зокрема визначення умов, за яких резонансні включення не лише зменшують хвильову відповідь, а й знижують crack-driving force у крихкій матриці.

\subsection*{Симуляція 4. Бетон із вбудованим ауксетичним або метаматеріальним армуванням}

Буде побудовано модель бетонної матриці з вбудованими 3D-друкованими ауксетичними або метаматеріальними армувальними елементами. Такий підхід відповідає сучасній концепції використання архітектурного армування для перерозподілу напружень, підвищення несучої здатності та керування траєкторіями тріщин у бетоні \cite{Vandadi2025}.

У межах цієї групи симуляцій буде розглянуто різні типи армувальної архітектури, зокрема bowtie-, tubular- або lattice-подібні структури. Особливу увагу буде приділено моделюванню контакту і міжфазної взаємодії між бетонною матрицею та армувальним елементом. Для опису можливого дебондингу буде використано когезійні закони змішаного руйнування. Передбачається аналіз сценаріїв, за яких тріщина проходить крізь армувальний елемент, відхиляється вздовж інтерфейсу, кiнкує назад у бетонну матрицю або зупиняється на архітектурному бар'єрі.

Очікуваний результат --- визначення умов ефективної роботи вбудованого архітектурного армування як внутрішньої ``пастки тріщин'' і встановлення вимог до геометрії армувальної структури та міцності міжфазного зчеплення.

\subsection*{Симуляція 5. Міжфазна деградація у бетонополімерних і цементополімерних метакомпозитах}

Буде розроблено моделі міжфазного руйнування у гібридних бетонополімерних і цементополімерних композитах, у яких полімерні або полімер-композитні елементи виконують функцію архітектурних вставок, решіток, покриттів, ремонтних шарів або елементів керування тріщиною. Такі системи є перспективними для створення багатофункціональних інфраструктурних матеріалів, однак їхня працездатність істотно залежить від міцності та довговічності меж полімер--цемент, армування--матриця або ремонтний шар--основа \cite{Zhang2025,Jipa2022}.

У симуляціях буде використано когезійні моделі змішаного руйнування, які дають змогу описати нормальне відшарування, зсувне руйнування та комбіновану деградацію інтерфейсу. Передбачається врахування впливу жорсткості полімерної фази, енергії руйнування міжфазної межі, початкових дефектів, орієнтації 3D-друку та можливого зниження механічних властивостей полімеру під дією лужного або вологого середовища. Окремо буде досліджено, як деградація інтерфейсу впливає на траєкторію тріщини, енергопоглинання та довговічність композиту.

Очікуваний результат --- побудова моделей, які дозволяють оцінити, за яких умов полімерні архітектурні елементи підвищують тріщиностійкість композиту, а за яких стають слабкою ланкою через міжфазне відшарування або середовищно-індуковане пошкодження.

\subsection*{Симуляція 6. Модельний фрагмент дорожнього або інфраструктурного шару з внутрішніми ``пастками тріщин''}

На завершальному етапі чисельної частини буде побудовано модельний фрагмент дорожнього, плитного або захисного інфраструктурного шару з однією або кількома внутрішніми архітектурними зонами, призначеними для керування траєкторією тріщини. Такими зонами можуть бути ауксетичні вставки, полімерні решітки, слабкі або контрольовано міцні міжфазні прошарки, резонансні включення або комбінації цих елементів.

Передбачається розгляд задачі поширення тріщини з початкового дефекту або концентратора напружень у присутності архітектурних елементів. Буде досліджено, чи може внутрішня архітектура змінити напрям тріщини, зменшити швидкість її розвитку, розподілити пошкодження в ширшій зоні або зупинити тріщину до досягнення критичного стану. Для цього будуть використані когезійно-зонні моделі, енергетичні критерії руйнування та параметричний аналіз геометрії архітектурних елементів.

Очікуваний результат --- формування узагальнених рекомендацій щодо проєктування внутрішніх архітектурних елементів, які можуть працювати як ``пастки тріщин'' у метаматеріально-орієнтованих інфраструктурних композитах з акцентом на системи дорожнього покриття.

Узагальнено, заплановані чисельні симуляції мають забезпечити перехід від аналізу окремих unit-cell моделей до оцінювання поведінки модельних фрагментів інфраструктурних шарів. Це дасть змогу встановити, як внутрішня архітектура композиту впливає на параметри руйнування, і сформулювати науково обґрунтовані принципи проєктування тріщиностійких метаматеріально-орієнтованих композитів для дорожньої та іншої цивільної інфраструктури.


% ------------------------------------------------------------
\section{Структура досліджень та етапи виконання роботи}
% Для пункту 15.4 Запиту і пункту 8 Технічного завдання.
% Дворічна структура.

Виконання роботи передбачається у два взаємопов'язані етапи. Перший етап буде спрямований на формування теоретичної та чисельної основи дослідження: вибір характерних типів метаматеріально-орієнтованих інфраструктурних композитів, побудову базових скінченноелементних і когезійно-зонних моделей, визначення основних механізмів зародження та поширення тріщин у матеріалах з внутрішньою архітектурою. Другий етап буде присвячений розвитку моделей для статичного й динамічного навантаження, урахуванню міжфазної деградації та середовищних чинників, чисельно-експериментальній верифікації результатів і формулюванню рекомендацій щодо створення тріщиностійких інфраструктурних композитів з акцентом на системи дорожнього покриття.

\subsection*{Етап 1. Розроблення моделей руйнування та вибір базових архітектур метаматеріально-орієнтованих інфраструктурних композитів}

На першому етапі передбачається провести систематизацію сучасних підходів до створення метаматеріально-орієнтованих композитів інфраструктурного призначення, включаючи метабетон із резонансними включеннями, ауксетичні цементні та бетонні структури, бетон із вбудованим архітектурним армуванням, багатофункціональні бетонополімерні системи та композити з вираженими міжфазними межами. На основі цього аналізу буде обрано кілька характерних типів внутрішньої архітектури, перспективних з погляду керування руйнуванням: комірчасті та ауксетичні структури, перфоровані цементні елементи, резонансні включення, полімерні або полімер-композитні армувальні решітки, а також шаруваті й міжфазні системи.

Для вибраних архітектур буде побудовано скінченноелементні моделі напружено-деформованого стану за статичного навантаження. Особливу увагу буде приділено зонам концентрації напружень, локалізації деформацій, шарнірним зонам ауксетичних комірок, лігаментам, межам включення--матриця та полімер--цементним інтерфейсам. На цьому етапі також передбачається побудова базових когезійно-зонних моделей для опису зародження тріщин, їхнього поширення в матриці, руйнування міжфазних меж, crack-bridging механізмів і посттріщинної поведінки цементних або бетонних композитів.

У межах першого етапу буде виконано параметричний аналіз впливу геометрії внутрішньої архітектури на умови зародження тріщин і початкову стадію їхнього поширення. Будуть досліджені характерні сценарії взаємодії тріщини з архітектурними елементами: проходження через лігамент, відхилення траєкторії на межі поділу фаз, поширення вздовж інтерфейсу, кiнкування в матрицю та зупинка на архітектурному бар'єрі. Це дасть змогу визначити початкові критерії ефективності внутрішніх елементів, які можуть виконувати функцію ``пасток тріщин''.

Очікуваним результатом першого етапу є створення базового комплексу механіко-математичних і скінченноелементних моделей руйнування метаматеріально-орієнтованих інфраструктурних композитів за статичного навантаження, вибір репрезентативних архітектур для подальшого аналізу, а також формування початкових чисельних карт локалізації напружень, зон зародження тріщин і можливих траєкторій їхнього поширення.

\subsection*{Етап 2. Моделювання статичного і динамічного руйнування, міжфазної деградації та експериментальна верифікація}

На другому етапі передбачається розвиток побудованих моделей для аналізу складніших умов навантаження, характерних для дорожньої та іншої цивільної інфраструктури. Буде досліджено поведінку метаматеріально-орієнтованих композитів за динамічного, імпульсного та циклічного навантаження, зокрема для систем із резонансними включеннями, віброзахисними архітектурними елементами та ауксетичними комірчастими структурами. Основну увагу буде приділено зв'язку між динамічною функціональністю матеріалу, розподілом енергії, зменшенням локальних напружень і параметрами руйнування.

Окремий блок робіт буде присвячено моделюванню міжфазної деградації у гібридних бетонополімерних і цементополімерних метакомпозитах. Будуть розглянуті межі полімер--цемент, армування--матриця, включення--покриття--матриця та ремонтний шар--основа. Для таких систем планується побудова когезійних моделей змішаного руйнування, які дадуть змогу описати дебондинг, відхилення тріщини вздовж межі поділу, її проникнення в матрицю або архітектурний елемент, а також втрату зчеплення внаслідок накопичення пошкоджень. Додатково буде проаналізовано вплив властивостей полімерних і полімер-композитних елементів, орієнтації 3D-друку та можливого середовищно-індукованого пошкодження на довговічність композиту.

Передбачається також чисельно-експериментальна верифікація запропонованих моделей. Для цього будуть використані випробування на стиск, згин, непрямий розтяг і SCB-випробування, які дають змогу визначати параметри тріщиностійкості, енергопоглинання, посттріщинної поведінки та міжфазного руйнування. Для виготовлення окремих архітектурних елементів, вставок, решіток або модельних зразків передбачається використання технологій 3D-друку полімерних чи полімер-композитних матеріалів. Отримані експериментальні дані будуть використані для ідентифікації параметрів когезійних моделей, перевірки чисельних карт руйнування та оцінки ефективності запропонованих архітектурних рішень.

Очікуваним результатом другого етапу є розроблення уточнених моделей руйнування метаматеріально-орієнтованих інфраструктурних композитів за статичного, циклічного й динамічного навантаження, побудова чисельних карт руйнування для вибраних типів внутрішньої архітектури, визначення параметрів міжфазної деградації та формулювання рекомендацій щодо створення тріщиностійких композитів для систем дорожнього покриття, плитних елементів, мостових настилів, захисних і ремонтних шарів.

Узагальненим результатом виконання двох етапів буде створення науково обґрунтованого підходу до аналізу та проєктування метаматеріально-орієнтованих інфраструктурних композитів, у яких внутрішня архітектура використовується не лише для керування жорсткістю, деформаційною здатністю чи динамічною відповіддю, а й для підвищення опору зародженню та поширенню тріщин.


% ------------------------------------------------------------
\section{Матеріально-технічна база}
% Для пункту 15.5 Запиту.

Для виконання проєкту передбачається використання наявної наукової та обчислювальної бази Інституту механіки ім. С.П. Тимошенка НАН України, а також спеціалізованого обладнання для виготовлення модельних зразків, проведення механічних випробувань і чисельно-експериментальної верифікації моделей руйнування метаматеріально-орієнтованих інфраструктурних композитів. Поєднання чисельного моделювання, 3D-друку архітектурних елементів та експериментального визначення параметрів тріщиностійкості дасть змогу реалізувати повний цикл досліджень: від побудови розрахункових моделей і вибору внутрішньої архітектури композиту до виготовлення прототипів, їх випробування та ідентифікації параметрів когезійних моделей.

Основним обладнанням для механічних випробувань має стати універсальна випробувальна машина Matest UNIMEC 300 з максимальним навантаженням 300 кН. Її використання дасть змогу проводити випробування цементних, бетонних, полімерцементних і модельних композитних зразків за схемами стиску, розтягу та згину. Передбачене використання захватів для круглих і плоских зразків, плит для випробувань на стиск, пристрою для трьох- і чотириточкового згину, а також програмного забезпечення SmartLab для реєстрації, обробки та аналізу експериментальних даних. Це обладнання є необхідним для визначення жорсткості, міцності, енергопоглинання, посттріщинної поведінки та граничних станів метаматеріально-орієнтованих композитів за квазістатичного навантаження.

Для спеціалізованих випробувань дорожніх, цементних і бетонних матеріалів передбачається використання універсального багатофункціонального тестера S205M Unitronic з відповідними пристроями та аксесуарами. Зокрема, планується застосування базового пристрою для випробувань на непрямий розтяг циліндричних зразків, пристрою для SCB-випробувань, навантажувального поршня, лінійного тензометричного перетворювача, калібрувальних засобів і вбудованого програмного забезпечення для SCB-випробувань. Такі випробування є особливо важливими для задач механіки руйнування дорожніх та інфраструктурних композитів, оскільки дають змогу визначати параметри тріщиностійкості, енергію руйнування, характер зародження і поширення тріщин, а також ефективність міжфазного зчеплення у гібридних матеріалах.

Для виготовлення модельних зразків, архітектурних вставок, армувальних решіток, комірчастих структур, полімерних або полімер-композитних елементів передбачається використання 3D-принтера CreatBot PEEK-250. Це обладнання дасть змогу створювати прототипи внутрішніх архітектур метаматеріально-орієнтованих композитів, зокрема ауксетичні решітки, перфоровані елементи, вставки для керування траєкторією тріщини, модельні ``пастки тріщин'', а також форми або допоміжні елементи для виготовлення бетонополімерних і цементополімерних зразків. Можливість 3D-друку є принципово важливою для цього проєкту, оскільки внутрішня геометрія композиту розглядається як один із головних параметрів керування руйнуванням.

Чисельна частина досліджень виконуватиметься із застосуванням методів скінченних елементів, когезійно-зонного моделювання та енергетичних критеріїв руйнування. Передбачається використання наявних програмних напрацювань колективу з моделювання тріщин, когезійних зон і міжфазного руйнування, а також розроблення нових розрахункових алгоритмів для метаматеріально-орієнтованих інфраструктурних композитів. Чисельні моделі будуть застосовані для аналізу напружено-деформованого стану, локалізації пошкоджень, взаємодії тріщин з архітектурними елементами, динамічного навантаження, міжфазної деградації та впливу геометричних параметрів внутрішньої структури на тріщиностійкість.

Зазначена матеріально-технічна база забезпечить можливість виконання повного комплексу робіт за темою: створення розрахункових моделей метаматеріально-орієнтованих композитів, виготовлення модельних архітектурних елементів, проведення випробувань на стиск, згин, непрямий розтяг і SCB, ідентифікація параметрів когезійних моделей, побудова чисельних карт руйнування та формулювання рекомендацій щодо проєктування тріщиностійких інфраструктурних композитів з акцентом на системи дорожнього покриття.


% ------------------------------------------------------------
\section{Досвід і доробок авторів}
% Для пункту 15.3 Запиту.
% Тут треба пов'язати тему з:
% - механікою руйнування;
% - CZM;
% - тріщинами в композитах;
% - чисельним моделюванням;
% - метаматеріалами;
% - експериментальною базою.

% ------------------------------------------------------------
\section{Очікувані наукові та науково-практичні результати}
% Для пункту 18.1 Запиту і пункту 10 Технічного завдання.

\begin{enumerate}
	\item Механіко-математична модель руйнування метаматеріально-орієнтованих інфраструктурних композитів.
	\item Алгоритми визначення умов зародження, поширення, відхилення та зупинки тріщин.
	\item Когезійні моделі міжфазної деградації в системах полімер--цемент, армування--матриця та soft layer--mortar.
	\item Чисельні карти руйнування для базових архітектур метаматеріально-орієнтованих композитів.
	\item Рекомендації щодо геометричних параметрів архітектурних елементів, які можуть виконувати функцію ``пасток тріщин''.
	\item Методика експериментальної верифікації з використанням випробувань на стиск, згин, непрямий розтяг і SCB.
	\item Пропозиції щодо застосування результатів у системах дорожнього покриття, захисних шарах, плитах, мостових настилах і ремонтних елементах.
\end{enumerate}

% ------------------------------------------------------------
\section{Наукова новизна та рівень розробки}
% Для пункту 17 Запиту.

Наукова новизна роботи полягає у формуванні fracture-mechanics підходу до аналізу та проєктування метаматеріально-орієнтованих інфраструктурних композитів, у яких внутрішня архітектура розглядається не лише як засіб керування жорсткістю, деформаційною здатністю, ефективним коефіцієнтом Пуассона, віброізоляційними властивостями або енергопоглинанням, а й як активний чинник впливу на зародження, поширення, відхилення та зупинку тріщин.

На відміну від наявних досліджень, у яких метабетон, ауксетичні цементні структури, бетон із вбудованим архітектурним армуванням або багатофункціональні бетонополімерні системи зазвичай розглядаються з позицій окремих ефектів, у цій роботі пропонується цілісний підхід до аналізу їхнього руйнування. Він поєднує опис внутрішньої архітектури композиту, локалізації напружень, crack-bridging механізмів, посттріщинної поведінки матриці, міжфазної деградації, динамічного навантаження та впливу експлуатаційного середовища.

Новими науковими результатами, які передбачається отримати, є:

\begin{enumerate}
	\item Механіко-математичні моделі руйнування метаматеріально-орієнтованих інфраструктурних композитів, що враховують вплив внутрішньої комірчастої, ауксетичної, резонансної, перфорованої або армувальної архітектури на умови зародження і поширення тріщин.

	\item Когезійно-зонні моделі руйнування характерних елементів архітектурних композитів: шарнірних зон ауксетичних комірок, лігаментів, меж включення--матриця, полімер--цементних інтерфейсів, армування--матриця та ремонтний шар--основа.
	
	\item Критерії оцінювання ефективності внутрішніх архітектурних елементів як ``пасток тріщин'', що можуть сприяти відхиленню, кiнкуванню, зупинці або стабілізації поширення тріщини.
	
	\item Чисельні карти руйнування для характерних типів метаматеріально-орієнтованих інфраструктурних композитів, які пов'язують геометричні параметри внутрішньої архітектури, властивості матриці та міжфазних меж із тріщиностійкістю, енергопоглинанням і посттріщинною поведінкою.
	
	\item Підходи до оцінювання впливу динамічної функціональності метабетону, резонансних включень і віброзахисних архітектурних елементів на локальні параметри руйнування, зокрема на максимальні розтягувальні напруження, crack-driving force та умови динамічного поширення тріщин.
	
	\item Моделі міжфазної деградації та середовищно-індукованого пошкодження у бетонополімерних і цементополімерних метакомпозитах, що враховують можливе відшарування, зсувне руйнування, втрату зчеплення та зміну властивостей полімерних або полімер-композитних архітектурних елементів.
	
\end{enumerate}

Рівень розробки запропонованого дослідження можна охарактеризувати як такий, що відповідає сучасним світовим тенденціям у галузі механічних метаматеріалів, архітектурних цементних матеріалів, метабетону, ауксетичних структур і багатофункціональних інфраструктурних композитів. Водночас у світовій літературі ці напрями переважно розвиваються окремо: досліджуються хвильові ефекти, ауксетична поведінка, енергопоглинання, самодіагностика або підвищення міцності, тоді як системний опис руйнування таких матеріалів залишається недостатньо сформованим.

Запропонована робота відрізняється від наявних досліджень тим, що переносить центр уваги з демонстрації окремих метаматеріальних ефектів на механіку тріщиноутворення, міжфазне руйнування та довговічність матеріалів інфраструктурного призначення. Такий підхід є особливо актуальним для систем дорожнього покриття, мостових настилів, плитних елементів, захисних і ремонтних шарів, де саме зародження та розвиток тріщин часто визначають експлуатаційну надійність і строк служби конструкції.

За власною оцінкою, робота не має прямих аналогів в Україні, оскільки поєднує механіку руйнування, когезійно-зонне моделювання, метаматеріально-орієнтоване проєктування внутрішньої архітектури та задачі довговічності дорожніх і цивільних інфраструктурних композитів. На міжнародному рівні робота відповідає актуальному напряму розвитку architected infrastructure materials, але має власний фокус --- створення моделей і критеріїв, які дають змогу використовувати внутрішню архітектуру композиту як інструмент керування тріщинами та підвищення опору руйнуванню.


% ------------------------------------------------------------
\section{Шляхи використання результатів}
% Для пунктів 18.2 і 18.3 Запиту.

Результати роботи матимуть наукове, методичне та прикладне значення для розвитку механіки руйнування композитних матеріалів інфраструктурного призначення. Розроблені механіко-математичні моделі, когезійно-зонні підходи, чисельні карти руйнування та критерії ефективності внутрішніх ``пасток тріщин'' можуть бути використані для аналізу, проєктування й оптимізації метаматеріально-орієнтованих композитів, призначених для роботи за статичного, циклічного та динамічного навантаження.

\subsection{Подальше використання у суспільній практиці}

Отримані результати можуть бути використані при створенні нових тріщиностійких композитних матеріалів для дорожньої, транспортної та цивільної інфраструктури. Зокрема, розроблені моделі та рекомендації можуть бути застосовані для обґрунтування внутрішньої архітектури дорожніх покриттів, плитних елементів, мостових настилів, захисних і ремонтних шарів, а також бетонополімерних і цементополімерних систем, у яких важливими є тріщиностійкість, енергопоглинання, міжфазна міцність і довговічність.

Практичне використання результатів передбачається у вигляді рекомендацій щодо вибору геометричних, матеріальних і міжфазних параметрів метаматеріально-орієнтованих інфраструктурних композитів. Такі рекомендації можуть стосуватися розміщення архітектурних елементів у матеріалі, вибору типу комірчастої або ауксетичної структури, параметрів полімерних або полімер-композитних армувальних вставок, характеристик міжфазного зчеплення та умов, за яких внутрішня архітектура сприяє відхиленню, зупинці або стабілізації поширення тріщин.

Розроблені чисельні моделі можуть бути використані як інструмент попереднього проєктування і порівняльного аналізу різних архітектурних рішень без необхідності виготовлення великої кількості експериментальних зразків. Це дасть змогу зменшити витрати на пошук оптимальної структури матеріалу, скоротити обсяг пробних випробувань і підвищити обґрунтованість вибору композитних систем для конкретних умов експлуатації.

Результати дослідження також можуть бути використані у науково-освітній діяльності, підготовці публікацій у фахових міжнародних і вітчизняних виданнях, доповідей на наукових конференціях, а також при розробленні нових наукових тем, дисертаційних досліджень і навчально-методичних матеріалів з механіки руйнування, механіки композитів, чисельного моделювання та інженерії матеріалів.

\subsection{Потенційні споживачі результатів}

Потенційними споживачами результатів роботи можуть бути науково-дослідні установи, заклади вищої освіти, проєктні організації та підприємства, діяльність яких пов'язана з розробленням, випробуванням, експлуатацією або ремонтом дорожніх, транспортних і цивільних інфраструктурних об'єктів.

До кола потенційних споживачів належать організації дорожньої галузі, що займаються проєктуванням, будівництвом, ремонтом і моніторингом дорожніх покриттів; установи та підприємства, пов'язані з експлуатацією мостових настилів, плитних конструкцій, аеродромних покриттів, захисних споруд і транспортної інфраструктури; виробники цементних, бетонних, полімерцементних, бетонополімерних і композитних матеріалів; лабораторії з механічних випробувань будівельних і дорожніх матеріалів; а також організації, зацікавлені у впровадженні 3D-друкованих архітектурних елементів у композитні матеріали інфраструктурного призначення.

Окрему групу потенційних споживачів становлять наукові колективи, які працюють у галузях механіки руйнування, механіки композитних матеріалів, будівельного матеріалознавства, механічних метаматеріалів, чисельного моделювання та структурної оптимізації. Для таких колективів результати роботи можуть бути корисними як основа для подальшого розвитку моделей руйнування матеріалів з внутрішньою архітектурою, дослідження міжфазної деградації та створення нових підходів до керування тріщиноутворенням у складних композитних системах.

У перспективі результати роботи можуть бути використані при формуванні науково обґрунтованих рекомендацій для створення тріщиностійких матеріалів дорожнього та інфраструктурного призначення, а також для підготовки пропозицій щодо подальших експериментальних, дослідно-конструкторських або прикладних робіт у сфері довговічних композитних матеріалів для цивільної інфраструктури.

% ------------------------------------------------------------
\section{Об'єкти права інтелектуальної власності}
% Для пункту 19 Запиту.
% Варіанти:
% - якщо наявних ОІВ не використовуємо: "не передбачається";
% - якщо плануємо створення: алгоритм, методика, програмний модуль,
% конструктивна схема архітектурного композиту.

% ------------------------------------------------------------
\section{Висновок}

Запропонована робота спрямована на розвиток механіки руйнування метаматеріально-орієнтованих інфраструктурних композитів, у яких внутрішня архітектура матеріалу розглядається як активний інструмент керування зародженням, поширенням, відхиленням і зупинкою тріщин. Такий підхід дає змогу перейти від традиційного удосконалення складу матеріалу до цілеспрямованого проєктування його внутрішньої структури з урахуванням тріщиностійкості, енергопоглинання, посттріщинної поведінки, міжфазної деградації та довговічності.

Особливість роботи полягає у поєднанні ідей механічних метаматеріалів, архітектурних цементних і бетонних матеріалів, ауксетичних структур, метабетону, 3D-друкованого армування та когезійно-зонного моделювання в межах єдиного fracture-mechanics підходу. Основний акцент зроблено на системах дорожнього покриття, однак отримані результати можуть бути поширені також на мостові настили, плитні елементи, захисні й ремонтні шари та інші об'єкти цивільної інфраструктури, для яких критичними є тріщиностійкість і довготривала експлуатаційна надійність.

У результаті виконання роботи очікується створення механіко-математичних моделей, чисельних методик, критеріїв оцінювання та рекомендацій щодо проєктування метаматеріально-орієнтованих інфраструктурних композитів із підвищеним опором руйнуванню. Розроблені підходи дадуть змогу оцінювати ефективність внутрішніх архітектурних елементів як ``пасток тріщин'', прогнозувати розвиток пошкоджень за статичного й динамічного навантаження та обґрунтовувати вибір геометричних, матеріальних і міжфазних параметрів перспективних композитних систем.

Таким чином, запропоноване дослідження має як фундаментальне значення для розвитку механіки руйнування матеріалів з внутрішньою архітектурою, так і прикладне значення для створення довговічних, тріщиностійких і функціонально адаптованих композитів дорожнього та інфраструктурного призначення.


% ------------------------------------------------------------

\begin{thebibliography}{99}
	
	\bibitem{Moini2024}
	Moini, R.
	Perspectives in architected infrastructure materials.
	\textit{RILEM Technical Letters}, 8, 125--140, 2023.
	\url{https://doi.org/10.21809/rilemtechlett.2023.183}
	
	\bibitem{Premadasa2026}
	Premadasa, R., Rouhbakhsh, S., Ghimire, S., Almasi, P., Kang, H., Wan, Z., Jiao, P., Zhang, Q.
	Emerging mechanical metamaterials for future civil engineering.
	\textit{Materials Today}, 98, 103419, 2026.
	\url{https://doi.org/10.1016/j.mattod.2026.103419}
	
	\bibitem{Wang2025}
	Wang, H., Zhao, S., Xu, C., Sun, K., Fan, R.
	Engineering metamaterials for civil infrastructure: From acoustic performance to programmable mechanical responses.
	\textit{Materials}, 18(17), 4032, 2025.
	\url{https://doi.org/10.3390/ma18174032}
	
	\bibitem{Alavi2026}
	Alavi, A. H.
	Metamaterials break the constraints of traditional civil infrastructure.
	\textit{Advanced Materials}, 38(6), e18198, 2026.
	\url{https://doi.org/10.1002/adma.202518198}
	
	\bibitem{Mitchell2014}
	Mitchell, S. J., Pandolfi, A., Ortiz, M.
	Metaconcrete: Designed aggregates to enhance dynamic performance.
	\textit{Journal of the Mechanics and Physics of Solids}, 65, 69--81, 2014.
	\url{https://doi.org/10.1016/j.jmps.2014.01.003}
	
	\bibitem{Imagawa2025}
	Imagawa, K., Ohno, M., Koga, Y., Ishida, T.
	Cement-based mechanical metamaterials with spiral resonators for vibration control.
	\textit{Cement and Concrete Composites}, 163, 106191, 2025.
	\url{https://doi.org/10.1016/j.cemconcomp.2025.106191}
	
	\bibitem{Xu2020}
	Xu, Y., Zhang, H., Schlangen, E., Lukovi{\'c}, M., {\v{S}}avija, B.
	Cementitious cellular composites with auxetic behavior.
	\textit{Cement and Concrete Composites}, 111, 103624, 2020.
	\url{https://doi.org/10.1016/j.cemconcomp.2020.103624}
	
	\bibitem{Xie2025}
	Xie, J., He, S., Xu, Y., Meng, Z., Zhou, W., Schlangen, E., {\v{S}}avija, B.
	Enhanced elastomer-like auxetic cementitious materials through strain-hardening cementitious composites (SHCC) with extended softening properties.
	\textit{Cement and Concrete Composites}, 161, 106069, 2025.
	\url{https://doi.org/10.1016/j.cemconcomp.2025.106069}
	
	\bibitem{Chen2024}
	Chen, M., Fang, S., Wang, G., Xuan, Y., Gao, D., Zhang, M.
	Compressive and flexural behaviour of engineered cementitious composites based auxetic structures: An experimental and numerical study.
	\textit{Journal of Building Engineering}, 86, 108999, 2024.
	\url{https://doi.org/10.1016/j.jobe.2024.108999}
	
	\bibitem{Vo2025}
	Vo, T. S., Kim, D. J.
	Auxetic meta-concrete with customized materials and structures: Experiments and simulations.
	\textit{Journal of Building Engineering}, 114, 114425, 2025.
	\url{https://doi.org/10.1016/j.jobe.2025.114425}
	
	\bibitem{Vandadi2025}
	Vandadi, M., Heidarnezhad, S., Pourhaji, P., Rahbar, N.
	Integrating 3D-printed auxetic structures for advanced concrete reinforcement.
	\textit{Advanced Materials Interfaces}, 12(17), e00095, 2025.
	\url{https://doi.org/10.1002/admi.202500095}
	
	\bibitem{Dong2023}
	Dong, S., Zhang, W., Wang, X., Han, B.
	New-generation pavement empowered by smart and multifunctional concretes: A review.
	\textit{Construction and Building Materials}, 402, 132980, 2023.
	\url{https://doi.org/10.1016/j.conbuildmat.2023.132980}
	
	\bibitem{Qin2024}
	Qin, H., Ding, S., Ashour, A., Zheng, Q., Han, B.
	Revolutionizing infrastructure: The evolving landscape of electricity-based multifunctional concrete from concept to practice.
	\textit{Progress in Materials Science}, 145, 101310, 2024.
	\url{https://doi.org/10.1016/j.pmatsci.2024.101310}
	
	\bibitem{Barri2023}
	Barri, K., Zhang, Q., Kline, J., Lu, W., Luo, J., Sun, Z., Taylor, B. E., Sachs, S. G., Khazanovich, L., Wang, Z. L., Alavi, A. H.
	Multifunctional nanogenerator-integrated metamaterial concrete systems for smart civil infrastructure.
	\textit{Advanced Materials}, 35(14), 2211027, 2023.
	\url{https://doi.org/10.1002/adma.202211027}
	
	\bibitem{Zhang2025}
	Zhang, B., Qiu, W., Niu, Y., Wang, B., Xu, S., Chen, J., Zhong, Y.
	Study on low-temperature shear bond performance of polymer-concrete interface.
	\textit{Construction and Building Materials}, 503, 144534, 2025.
	\url{https://doi.org/10.1016/j.conbuildmat.2025.144534}
	
	\bibitem{Jipa2022}
	Jipa, A., Reiter, L., Flatt, R. J., Dillenburger, B.
	Environmental stress cracking of 3D-printed polymers exposed to concrete.
	\textit{Additive Manufacturing}, 58, 103026, 2022.
	\url{https://doi.org/10.1016/j.addma.2022.103026}
	
\end{thebibliography}


\end{document}